\subsection{拾球控制状态机}\label{sec:fsm}

机器人要把场上的网球逐个拾起再倒进收集桶，这是一连串有先后次序的动作，而不是一个可以一次算完的
函数。它得先找到球、追上去、把球夹住，再转身找到桶、靠近桶、把球松开，然后回到找球。每一步成立
的前提都依赖上一步的结果，且依赖此刻相机看到了什么。把这套行为组织成一台有限状态机，机器人在任一
时刻只处于一个明确的阶段，每帧依据当前阶段和最新的一次检测决定下一步怎么走，再在条件满足时切换到
下一个阶段。

整套流程由 \texttt{GameState} 枚举的五个状态构成，分别是追球（\texttt{CHASE\_BALL}）、夹取
（\texttt{GRAB}）、寻桶（\texttt{FIND\_BUCKET}）、近桶（\texttt{APPROACH\_BUCKET}）与投放
（\texttt{DEPOSIT}），机器人启动即进入追球。状态机自身不做任何输入输出，也不睡眠等待，它只接收
一个 \texttt{Detection} 并返回一个 \texttt{ControlOutput}，后者描述这一帧底盘该如何动
（前进、刹车或待机，附带左右轮速度）以及机械臂该做什么动作。这样一帧从相机到指令的时延就等于纯
计算时延，不会被状态机内部的忙等拖长。\textbf{这一点与参考实现不同，参考实现把夹取与投放写成在循环里忙等
数百毫秒的阻塞段，本实现把夹取与投放也建模成短暂的、非阻塞的显式状态，靠帧计数自然退出。}

\repfig{tennis-state-machine.png}{拾球控制状态机的五个状态与触发各次切换的条件，回路在投放完成后
回到追球，开始下一轮}{fig:fsm}

\paragraph{状态之间的流转}
追球状态把球的检测结果转成对底盘的转向与前进量。球的水平偏移决定转向偏置，球在画面里占的面积比
（\texttt{area\_ratio}）作为距离的代理量决定前进速度，远则全速、近则收速。当球足够近且持续若干帧
（\texttt{stop\_confirm\_cnt}）稳定地对准夹爪的目标位置时，状态机刹车并切入夹取。夹爪并不在光轴
正中，而是装在偏右的位置，因此对准的目标是一个带偏置的点（\texttt{stop\_center\_offset}）而非画面
中心。若某一帧没看到球，机器人先朝球最后出现的方向转动搜寻，搜寻若干帧仍未找回，便改为左右往复
扫视，扫视方向每隔固定帧数翻转一次。进入夹取后，机械臂的夹取动作只在该状态的第一帧下发一次，再经
若干帧（\texttt{grab\_settle\_frames}）安定后转入寻桶。寻桶时底盘原地旋转，一旦连续若干帧
（\texttt{bucket\_confirm\_cnt}）都看到桶，便切入近桶；近桶时按桶的偏移转向、按桶的面积比收速逼近，
中途若桶丢失超过若干帧（\texttt{bucket\_lost\_frames}）则退回寻桶，桶在画面里占满到投放阈值
（\texttt{bucket\_area\_deposit}）时刹车并切入投放。投放时机械臂的松开动作同样只在第一帧下发，
安定若干帧（\texttt{deposit\_settle\_frames}）后机械臂复位，状态机回到追球，一轮完整的拾球至此闭合。

\paragraph{一帧只跑一个检测器}
两路感知喂给状态机，球由运行在 NPU 上的 YOLOv8 检测（\ref{sec:npu}），桶则由色调饱和度明度
（HSV）颜色检测识别。读 \texttt{BucketDetector} 可以确认桶的识别全程不碰神经网络，它把 RGB888
转为 HSV，用一个跨越色环两端的双区间红色掩膜，再以两遍并查集连通域标记取出像素数最大且超过最小
面积阈值的红色块，把它的中心与面积比作为桶的观测。两个检测器的代价并不相同，前者要占用加速器跑
一次推理，后者只在 CPU 上扫一遍像素，因此让两者每帧都跑会平白浪费一次推理。状态机据此对外暴露一个
\texttt{perception\_mode()}，在追球与夹取两个状态下它返回 \texttt{BALL}，在其余状态下返回
\texttt{BUCKET}，感知线程依据这个值决定下一帧只运行哪一个检测器。这样同一帧相机图像永远只被一路
处理，绝不会为了同时拿到球与桶而把同一帧解码或处理两遍，这一点与参考实现在寻桶相关状态下把同一张
JPEG 解码两次的做法正相反。

\paragraph{把检测结果转为动作}
状态机消费的是最新的一帧相机图像（\ref{sec:camera}），输出的则是底盘指令与机械臂动词。控制器
（\texttt{Controller}）拿到 \texttt{ControlOutput} 后，把 \texttt{Drive} 译为对左右轮速度的设置、
把 \texttt{Brake} 与 \texttt{Standby} 译为刹车与待机，交给差速底盘执行（\ref{sec:chassis}），把
\texttt{Grab}、\texttt{Release}、\texttt{Ready} 三个机械臂动词交给机械臂执行（\ref{sec:arm}）。
为避免在热路径上反复下发完全相同的电机指令，控制器只在指令相较上一帧发生变化时才真正写下去。
至此，相机的一帧经感知、状态机决策、再到底盘与机械臂，构成一条从感知到控制的闭合回路，机器人就在
这条回路上一颗接一颗地把球拾起。
